\documentclass[11pt,a4paper]{article}
\usepackage[utf8]{inputenc}
\usepackage{amsmath,amssymb,amsfonts}
\usepackage{geometry}
\usepackage{hyperref}

\geometry{margin=1in}

\title{\textbf{\Huge Directed Graph Feedback Loop Stability \& Memory Buffers}\\
\Large Real-Time DSP Safety Guarantees in Dynamic Modular Graphs}
\author{\textbf{Time Dilation DAW (v4.0) Engineering Group}}
\date{\today}

\begin{document}

\maketitle

\begin{abstract}
This document formulates the cycle detection algorithms and 1-block delay buffer history allocation mechanism ensuring $100\%$ real-time audio thread stability in \textbf{Time Dilation DAW}.
\end{abstract}

\section{Modular Audio Directed Graph Representation}
The modular audio processing network is modeled as a Directed Graph $G = (V, E)$, where $V$ represents processing nodes and $E \subseteq V \times V$ represents audio or control patch connections.

\section{Feedback Loop Detection Algorithm}
A feedback loop occurs when there exists a directed path $v_1 \to v_2 \to \dots \to v_k \to v_1$.
Depth-First Search (DFS) with node visitation states $\text{State} \in \{\text{UNVISITED}, \text{VISITING}, \text{VISITED}\}$ identifies cyclic dependencies in $O(|V| + |E|)$ time complexity.

\section{1-Block History Delay Buffer Solution}
For any connection $e = (u, v) \in E$ forming a feedback loop cycle:
\begin{equation}
x_{\text{dest}}[n] = y_{\text{src, prevBlock}}[n], \quad \forall n \in [0, K_{\text{block}}-1]
\end{equation}
where $y_{\text{src, prevBlock}}$ is the audio buffer produced during block $N-1$.
This guarantees zero buffer starvation, zero thread deadlocks, and absolute DSP stability without requiring zero-delay matrix solvers.

\end{document}
